%
% Front matter
%

\begin{frontmatter}

\title{
  T2K 3.010$\times10^{20}$-POT 3-Flavour Muon-Neutrino Disappearance Analysis
}

\vspace*{0.7cm}

\author[RAL]     {Costas Andreopoulos}
\author[RAL,OXF] {Thomas Dealtry}
\author[IFIC]    {Lorena Escudero S\'anchez}
\author[LANCS]   {Nick Grant}
\author[RAL,OXF] {and Alfons Weber}

\address[RAL]{Particle Physics Dept., Rutherford Appleton Laboratory, STFC \\
              Harwell Oxford Campus\\
              Oxfordshire OX11 0QX, UK}

\address[OXF]{Subdepartment of Particle Physics, University of Oxford \\
              Denys Wilkinson Building, Keble Road\\
              Oxfordshire OX1 3RH, UK}

\address[IFIC]{Instituto de Fisica Corpuscular\\
              CSIC and University of Valencia\\
              Valencia, Spain}

\address[LANCS]{Physics Dept., Lancaster University \\
              Lancaster, LA1 4YB, UK}

%
% Abstract
%

\begin{abstract}

We report the results of the {\em VaLOR} (Valencia-Lancaster-Oxford-Rutherford) 
3-flavour $\nu_{\mu}$-disappearance analysis on the combined Run 1+2+3 dataset.
This analysis is an update of the work presented in
Refs. \cite{OurRun1DisappAnalysis}, \cite{OurRun12DisappAnalysis} and \cite{OurRun12DisappAnalysis3Flavour}
with improved methods and an enlarged dataset.
The Run 1+2+3 dataset corresponds to an integrated J-PARC neutrino beam
exposure of 3.010$\times$10$^{20}$ POT.

This analysis predicts 204.75 $\pm$ 16.75 (syst)
1 $\mu$-like ring events in SuperK in the absence of any oscillation,
but only 58 were observed. The observed deficit has a strong energy dependence;
the ratio of observed to expected, under the no-oscillation hypothesis, is
$\sim$30\% $<$ 0.5 GeV, 
$\sim$10\% between 0.5 and 1 GeV and 
$\sim$75\% $>$ 1 GeV.

A $\nu_{\mu}$-disappearance analysis was performed 
in a framework of 3-flavour oscillations including matter effects in constant-density matter.
Two versions of the oscillation fit were performed:
one (the `2+N' fit) allows all 48 systematics considered in this analysis to float, 
while the other (the `2+8' fit) allows only the 8 most dominant systematics to float
(the dominant systematics are those identified as having an effect on the measured oscillation parameters
that is greater than 25\% of the expected statistical error, for at least one set of true oscillation parameters in the area of interest). 

The `2+N' 3-flavour $\nu_{\mu}$-disappearance fit to
the observed reconstructed energy spectrum of 1 $\mu$-like ring events of the combined Run1+2+3 dataset,
assuming the first octant ($\sin^22\theta_{23} \leq 0.5$),
with $\sin^{2}2\theta_{23}$ constrained to the physical region,
yields
$|\Delta m^{2}_{32}|$ = 2.441$\times$10$^{-3}$ $eV^{2}/c^{4}$  and
$\sin^{2}\theta_{23}$ = 0.514
as the best-fit values (p-value = 0.80).
For the '2+N' fits for $\sin^22\theta_{23}$, assumming both the first ($\sin^22\theta_{23} \leq 0.5$)
and second octant ($\sin^22\theta_{23} \geq 0.5$),
the best-fit values are almost identical.
All fits used 2012 PDG values for
$\sin^{2}2\theta_{13}$, $\sin^{2}2\theta_{12}$ and $\Delta m^{2}_{21}$ \cite{PDG2012},
and assumed $\delta_{CP}$=0 and the normal mass hierarchy.

Confidence regions obtained using both the constant-$\Delta\chi^2$ method and the 
method of Feldman and Cousins are shown for both the `2+N' and `2+8' fits.
% Using the marginalized $\Delta\chi^2$ distributions obtained from the fits,
% and the constant-$\Delta\chi^2$ method,
% the 90\% CL allowed region, for both `2+N' and `2+8' fits, assuming the first octant, can be summarized as
% 2.14$\times$10$^{-3}$ $eV^{2}/c^{4}$ $< |\Delta m^{2}_{32}| < $ 2.76$\times$10$^{-3}$ $eV^{2}/c^{4}$ and
% $\sin^{2}2\theta_{23} > $ 0.957.
% The 90\% CL allowed region, for the '2+N' fit, assuming the second octant, can be summarized as
% 2.??$\times$10$^{-3}$ $eV^{2}/c^{4}$ $< |\Delta m^{2}_{32}| < $ 2.??$\times$10$^{-3}$ $eV^{2}/c^{4}$ and
% $\sin^{2}2\theta_{23} > $ 0.???.
Using the results of the '1+N' fits, 
and the constant-$\Delta\chi^2$ method,
the 90\% CL allowed region, for the `2+N' fit for $\sin^2\theta_{23}$, can be summarized as
$2.19$$\times$10$^{-3}$ eV$^{2}$/c$^{4}$ $< |\Delta m^2_{32}| < 2.72$$\times$10$^{-3}$ eV$^{2}$/c$^{4}$ and
$0.397 < \sin^{2}\theta_{23} < 0.630$.
The corresponding limits for the '2+N' fit assuming the first octant, can be summarized as
$2.19$$\times$10$^{-3}$ eV$^{2}$/c$^{4}$ $< |\Delta m^2_{32}| < 2.72$$\times$10$^{-3}$ eV$^{2}$/c$^{4}$ and
$\sin^{2}2\theta_{23}$ $>$ 0.958.
The corresponding limits for the '2+N' fit assuming the second octant, can be summarized as
$2.19$$\times$10$^{-3}$ eV$^{2}$/c$^{4}$ $< |\Delta m^2_{32}| < 2.72$$\times$10$^{-3}$ eV$^{2}$/c$^{4}$ and
$\sin^{2}2\theta_{23}$ $>$ 0.931.

\end{abstract}

%
% Keywords
%

%\begin{keyword}
%\end{keyword}

\end{frontmatter}

\clearpage

\tableofcontents

\clearpage

%\listoffigures
%\listoftables


